\documentclass[11pt]{article}
\usepackage[margin=25mm]{geometry}
\usepackage{fontspec}
\setmainfont{Latin Modern Roman}
\usepackage{amsmath,amssymb,mathtools}
\usepackage{unicode-math}
\setmathfont{Latin Modern Math}
\usepackage{xurl,hyperref}
\hypersetup{colorlinks=true,urlcolor=blue}
\setcounter{secnumdepth}{0}
\setlength{\parindent}{0pt}
\setlength{\parskip}{5pt}
\setlength{\emergencystretch}{3em}
\newcommand{\sref}[2]{\hyperlink{source:#1:#2}{[S#2]}}
\newcommand{\cataloguescope}[1]{\paragraph{Recorded target.}#1}
\begin{document}
\section{Erdős-Szekeres with Visibility}\label{statement}

\subsection{Definitions and mathematical statement}

\paragraph{Shared notation.}$\mathbb N=\{1,2,\ldots\}$, and $\mathbb Z,\mathbb Q,\mathbb R,\mathbb C$ denote the integers, rationals, reals and complex numbers. Any other conventions are specified locally.


For a finite set $A\subset\mathbb R^2$, two distinct points $x,y\in A$ are visible relative to $A$ if the open line segment between them contains no point of $A$; equivalently, $[x,y]\cap A=\{x,y\}$. A subset is pairwise visible if each pair of its distinct members is visible relative to the entire ambient set $A$, including points outside the subset.
\paragraph{Statement recorded in the supplied source.}
For every pair of integers $k,\ell\geq2$, does there exist an integer $n_0(k,\ell)$ such that every finite set $A\subset\mathbb R^2$ with $|A|\geq n_0(k,\ell)$ contains either $k$ distinct collinear points or $\ell$ pairwise visible points? \sref{statement}{2}
\subsection{Short English statement}
Does every sufficiently large planar point set contain either many points on one line or a prescribed number of points with no other point blocking any connecting segment?
\subsection{Sources}
\begin{description}
\item[\hypertarget{source:statement:1}{S1}] ULAM.AI and the UnsolvedMath Contributors, \emph{UnsolvedMath: A Curated Collection of Open Mathematics Problems}, catalogue entry. CC BY 4.0. \url{https://huggingface.co/datasets/ulamai/UnsolvedMath}.
\item[\hypertarget{source:statement:2}{S2}] Zachary Abel, Brad Ballinger, Prosenjit Bose, Sébastien Collette, Vida Dujmović, Ferran Hurtado, Scott D. Kominers, Stefan Langerman, Attila Pór and David R. Wood, \emph{Every Large Point Set Contains Many Collinear Points or an Empty Pentagon}, Graphs and Combinatorics 27 (2011), 47--60, §1.4, Conjecture 4. \url{https://arxiv.org/abs/0904.0262}.
\item[\hypertarget{source:statement:3}{S3}] Ben Green, \emph{100 Open Problems}, maintained author list, Problem 70 and its comments, accessed 7 October 2026. \url{https://people.maths.ox.ac.uk/greenbj/papers/open-problems.pdf}.
\end{description}
\label{end:statement}
\end{document}
